By homogeneity divide all variables by $d$, temporarily making $d=1$. The second and third equations give
\[
b=\dfrac{c(31-9c)}{4(c+1)},\qquad a=\dfrac{9c+25}{95-9c}.
\]
Positivity implies $0<c<\dfrac{31}{9}$, so both denominators are nonzero. Substituting these formulas into the first equation and simplifying gives
\[
\dfrac{5c(3c-5)^2(117c^2-1103c-700)}{16(c+1)^2(9c-95)}=0.
\]
The last quadratic is negative throughout $0<c<\dfrac{31}{9}$, because $117c-1103<0$. Thus $c=\dfrac{5}{3}$, and then $a=\dfrac{1}{2},b=\dfrac{5}{2}$. These ratios also satisfy the fourth equation; scaling by six gives the positive integer solution $(3,15,10,6)$. Therefore
\[
\boxed{\dfrac{1}{5}+\dfrac{3}{2}+\dfrac{5}{3}+2=\dfrac{161}{30}}.
\]
